% Auxiliary geometric-complexity estimate; no cover-atlas or finite-generation conclusion.
\begin{proposition}[Proved: degree-linear size of punctured-cover resolutions]
\label{D:punctured-cover-linear-resolution}
Let $(S,s)$ be a normal algebraic surface germ over an algebraically
closed field of positive characteristic.  Fix a projective resolution
$Y\to S$ which is an isomorphism away from $s$, with nonempty exceptional
curves $E_i$, and a strictly anti-nef positive cycle
\[
 A=\sum_i a_iE_i,\qquad c_i=-A\cdot E_i>0.
\]
Put
\[
 a=\sum_i a_i,\qquad b=K_Y\cdot A+a,\qquad c=3b+a,
 \qquad \lambda=\max_i\frac{b}{a_ic_i}.
 \tag{D.LR.1}
\]
For any connected finite etale cover of $S\setminus\{s\}$ of
degree $n$, normalize $Y$ in its field, resolve, and contract all
smooth rational $(-1)$-curves which map to points of $Y$.
Let $f:Z\to Y$ be the resulting proper generically finite morphism.
The cover need not be Galois.  Then
\[
 \#\operatorname{Irr}\bigl((f^*A)_{\mathrm{red}}\bigr)
 \le n(b+a)+(95+93\lambda)nc.
 \tag{D.LR.2}
\]
In particular the number of exceptional components, the sum of their
arithmetic genera, and the sum of the negatives of their
self-intersections are all bounded linearly in $n$, with constants
depending only on the fixed germ, resolution, and cycle.
The same assertion holds for covers of the punctured strict
henselization, with the same constants.

This is an estimate for resolutions of covers unramified outside
the isolated point.  It does not assert an object atlas, a bound on
the number of covers, or finite generation of the local group.
\end{proposition}

\begin{proof}
Use the notation of Propositions~\ref{D:exceptional-degree-linear-core}
and~\ref{D:weighted-exceptional-chain-compression}, with no original
divisorial boundary.  Write
\[
 C=f^*A=\sum_Hm_HH,\qquad
 R=K_Z-f^*K_Y=\sum_H\delta_HH.
 \tag{D.LR.3}
\]
The second divisor is the divisor of the nonzero determinant of
the differential.  It is effective and supported on the exceptional
fibre, because the cover and the chosen resolution morphism are
etale away from that fibre.  All $m_H$ are positive integers.
At a component $F$ dominating $E_i$, its coefficients satisfy
\[
              m_F=a_i e_F,\qquad \delta_F=d_F,
\]
where $d_F$ is the different exponent of the corresponding
extension of discrete valuation rings, including its full
possibly inseparable residual extension.

\par\smallskip\noindent\textbf{A maximum principle for ramification coefficients.}\quad
Proposition~\ref{D:surface-uniform-different} gives
$d_F/e_F\le b/c_i$ for a Galois cover.  It gives the same separate
bound for every finite cover by taking its finite Galois closure:
in a tower of finite separable DVR extensions $L/M/K$, transitivity
of the different gives
\[
 \frac{d(L/K)}{e(L/K)}
 =\frac{d(M/K)}{e(M/K)}
  +\frac{d(L/M)}{e(L/M)e(M/K)}
 \ge\frac{d(M/K)}{e(M/K)}.
\]
Thus $\delta_F/m_F\le\lambda$ on every dominating component,
independently of the cover degree.

For a point-mapped component $Q$, the chosen contractions give
$K_Z\cdot Q\ge0$, while projection gives $C\cdot Q=0$.
Writing $b_Q=-Q^2>0$, these statements become
\[
 \sum_{H\ne Q}(H\cdot Q)m_H=b_Qm_Q,
 \qquad
 \sum_{H\ne Q}(H\cdot Q)\delta_H\ge b_Q\delta_Q.
\]
Consequently
\[
 \sum_{H\ne Q}(H\cdot Q)m_H
       \left(\frac{\delta_H}{m_H}-\frac{\delta_Q}{m_Q}\right)
                         \ge0.                         \tag{D.LR.4}
\]
Each connected component of the exceptional graph contains a
dominating component.  Otherwise $C$ would have intersection zero
with every component of that connected component, contrary to
negative definiteness of its intersection matrix and the positive
coefficients of $C$.
If the maximum of the finitely many ratios $\delta_H/m_H$ were
larger than $\lambda$, it would occur at a point-mapped vertex.
Every summand in (D.LR.4) would then be nonpositive.  Equality
would force the same maximum at every adjacent vertex, and hence,
by connectedness, at a dominating vertex.  This is impossible.
We have proved the uniform estimate
\[
                         0\le\delta_H\le\lambda m_H
                         \quad\hbox{for every }H.        \tag{D.LR.5}
\]

\par\smallskip\noindent\textbf{The Jacobian cost of a long constant chain.}\quad
Consider one of the constant-multiplicity segments in
Proposition~\ref{D:weighted-exceptional-chain-compression}, of
common weight $m$.  Suppose it has $N-1\ge1$ internal vertices.
These form a chain of smooth rational $(-2)$-curves, and all map
to the same point of $Y$.  Contract just these internal curves.
The contraction argument of
Proposition~\ref{D:contract-zero-cost-blocks} applies to this
negative-definite subchain: every positive cycle $T$ on it has
$K_Z\cdot T=0$ and $T^2$ a negative even integer, so
$p_a(T)=1+T^2/2\le0$.  It gives a normal scheme contraction
$\rho:Z\to\overline Z$, with an $A_{N-1}$ rational double point.
The morphism $f$ descends to $\overline f:\overline Z\to Y$.
Indeed the graph image over $\overline Z$ is proper and
quasi-finite: the contracted chain has one image in $Y$ and all
other fibres are singletons.  The graph projection is therefore
finite birational onto the normal surface $\overline Z$, and
hence an isomorphism.

The completed local ring of an $A_{N-1}$ point in arbitrary
characteristic is
\[
                  B=k[[u,v,w]]/(uv-w^N)
\]
\citep[Theorem 0.4.13, printed p.~72, together with
Proposition 0.4.8]{C_Enriques2025}.
No assertion identifying the two retained end curves with the
coordinate branches is needed.  We use only the underlying
singularity and its minimal resolution.
For clarity, its standard resolution charts, indexed by
$i=0,\ldots,N-1$, are
\[
 u=x^{i+1}y^i,\qquad
 v=x^{N-i-1}y^{N-i},\qquad w=xy.
 \tag{D.LR.6}
\]
The intervening exceptional valuations $v_j$, $1\le j\le N-1$,
therefore satisfy
\[
                     v_j(u)=j,\qquad
                     v_j(v)=N-j,\qquad v_j(w)=1.        \tag{D.LR.7}
\]
These charts can be obtained successively by the normalized
blowup of $(u,w)$: its $u$-chart is smooth, with
$w=uz$ and $v=u^{N-1}z^N$, and its $w$-chart has
$u=wz$ and $zv=w^{N-1}$, so induction applies.
Adjacent charts in (D.LR.6) glue by
$x_{i+1}=y_i^{-1}$ and $y_{i+1}=x_i y_i^2$.
They give a proper resolution whose intervening curves have
normal degree $-2$.  These are the curves of the contracted
chain, since that chain is already its minimal resolution.

A generator of the dualizing module of $B$ is
\[
                         \eta=\frac{du\wedge dw}{u}.
\]
The expression is interpreted in the fraction field and extends
as a dualizing generator.  The hypersurface relation gives
\[
 du\wedge dw=u\eta,\qquad
 dv\wedge dw=-v\eta,\qquad
 du\wedge dv=Nw^{N-1}\eta.
\]
In particular the image of the module of continuous two-forms
in the dualizing module is
\[
                 (u,v,Nw^{N-1})\eta.                    \tag{D.LR.8}
\]
On each chart (D.LR.6), direct differentiation gives
$\rho^*\eta=dx\wedge dy$: the coefficient is
$(i+1)-i=1$, also in positive characteristic.  Thus the
resolution is crepant, with this explicit nonvanishing generator.

Choose regular parameters $z_1,z_2$ at the image point in $Y$.
The completed pullback of $dz_1\wedge dz_2$ is $h\eta$ for
some nonzero
\[
                         h\in(u,v,Nw^{N-1}).
\]
It is nonzero because the original cover is generically separable.
The coefficient $\delta_j$ of $R$ on the $j$th internal curve
is exactly $v_j(h)$, by the displayed crepant formula.  Since
coefficients in $B$ have nonnegative valuations, (D.LR.7)--(D.LR.8)
give, for $j=\lfloor N/2\rfloor$,
\[
 \delta_j\ge\min\{j,N-j,N-1\}=\lfloor N/2\rfloor.
 \tag{D.LR.9}
\]
If the characteristic divides $N$, the last generator of the
ideal is zero and the same inequality still follows from the
first two generators.  Combining (D.LR.9) with (D.LR.5) and
the constant multiplicity $m$ yields
\[
                         N-1\le2\delta_j\le2\lambda m.
                         \tag{D.LR.10}
\]
This argument is independent of the positions of the retained
end curves in the completed $A_{N-1}$ model.

\par\smallskip\noindent\textbf{Summing the chain lengths.}\quad
Let $P$ be the endpoint set in
Proposition~\ref{D:weighted-exceptional-chain-compression}.
There are at most $95nc$ actual noncore vertices after the
internal vertices of constant segments have been suppressed,
and $\sum_{Q\in P}m_Q\le31nc$.
Each endpoint lies on at most three segments.  Both endpoint
weights of a constant segment equal its weight $m$, so
\[
               \sum_{\text{constant segments}}m
                    \le\frac32\sum_{Q\in P}m_Q
                    \le\frac{93}{2}nc.
\]
By (D.LR.10), restoring all their internal vertices adds at
most $93\lambda nc$ vertices.  Adding the core bound
$n(b+a)$ proves (D.LR.2).
All noncore curves have arithmetic genus zero and square $-2$;
the core genus and self-intersection sums are bounded by
Proposition~\ref{D:exceptional-degree-linear-core}.  This proves
the remaining linear estimates.

For the strict-henselian assertion, each finite cover descends
to a finite-type etale neighborhood of the germ.  The chosen
exceptional curves, their intersection numbers, and the
residue field remain unchanged.  The preceding proof applies
there, exactly as in Proposition~\ref{D:surface-uniform-different}.
It produces the same constants.  No assertion about a
degree-linear presentation of the formal neighborhoods or
their covering algebras has been used.
\end{proof}
